%% =====================================================================
%% 卒業論文 LaTeX テンプレート
%%   玉川大学 工学部 情報通信工学科 計算機システム研究室
%%
%% ファイル: main.tex（本文），references.bib（参考文献），
%%           .latexmkrc（コンパイルの設定．変更しないこと）
%% Overleaf: Menu → Settings → Compiler を「LaTeX」にする
%% 手元・Docker: latexmk main.tex
%% =====================================================================
\documentclass[uplatex,dvipdfmx,a4paper,11pt,openany,nomag*]{jsreport}

% ---- 余白（上35mm，下・左・右30mm）----
\usepackage[top=35mm,bottom=30mm,left=30mm,right=30mm]{geometry}

% ---- よく使うパッケージ ----
\usepackage[haranoaji]{pxchfon}  % 日本語フォントをPDFに埋め込む
\usepackage{amsmath,amssymb}     % 数式
\usepackage{graphicx}            % 図の挿入
\usepackage{booktabs}            % 表の罫線
\usepackage{url}                 % URLの表示
\usepackage[hidelinks]{hyperref} % 目次などのリンク
\usepackage{pxjahyper}           % PDFのしおりの日本語化

% 参考文献の見出しを「参考文献」にする
\renewcommand{\bibname}{参考文献}

\begin{document}

%% =====================================================================
%% 表紙（年度・タイトル・学籍番号・氏名を書き換える）
%% =====================================================================
\begin{titlepage}
  \centering
  {\LARGE 20xx年度　卒業論文\par}
  \vspace{20mm}
  {\Huge タイトル\par}
  \vfill
  {\Large
    玉川大学　工学部　情報通信工学科\par
    計算機システム研究室\par
    学籍番号：xxxxxxxxx\par
    氏名：hogehoge\par}
  \vspace{60mm}
\end{titlepage}

%% =====================================================================
%% 概要
%%   \chapter* は番号なしの章．次の行で目次に追加している．
%% =====================================================================
\chapter*{概要}
\addcontentsline{toc}{chapter}{概要}

研究背景，問題提起，提案手法，この論文で実施したこと，達成したことを簡単にまとめる．
ここを読めばこの論文がどのようなことを目指して，どのようなことを実施した論文なのかわかるように書く．
記述量としてはここ1ページに収める．

（書き方のメモ）
章は \verb|\chapter{}|，節は \verb|\section{}|，項は \verb|\subsection{}| で作る．
目次はコンパイルすると自動で作られる．
段落の先頭は自動で字下げされるので全角スペースは入れない．
段落を変えるときは空行を1行入れる．

%% =====================================================================
%% 目次（自動で作られる）
%% =====================================================================
\tableofcontents

%% =====================================================================
%% 第1章 序論
%% =====================================================================
\chapter{序論}
\label{chap:intro}

\section{背景}

厚生労働省による調査では近年，少子高齢化が加速し〜

\section{研究目的}

\section{本論文の構成}

本論文は第\ref{chap:related}章でxxxxについて述べる．
第\ref{chap:method}章では提案手法について解説し，第\ref{chap:experiment}章では提案手法のxxxについてxxxするための実験を実施し，その結果を示す．
第\ref{chap:discussion}章では実験結果を踏まえ結果について考察し，今後の課題を示す．
第\ref{chap:conclusion}章で本論文の結論を述べる．

%% =====================================================================
%% 第2章 関連研究
%% =====================================================================
\chapter{関連研究}
\label{chap:related}

\section{XxxのFPGA実装}

関連する研究について参考文献を示し，その文献で実施されたこと，不足していることなどをレビューし，自分の研究の立場を明らかにする．
稲邑先生の“「卒業論文を書く」とはどういうことか？　これを知っておけば卒論を書くのが３倍楽になる”に書いてあるような先行研究を実施した人との会話になっているかどうか．
独り言になっていないか？

% ---- 引用の例（不要になったら消す）----
文献は \verb|\cite{キー}| で引用する（キーは references.bib に書いた名前）．
例えば，Zhangらはルーフラインモデルを用いてCNNアクセラレータの設計空間を探索した\cite{zhang2015fpga}．
複数の文献はまとめて引用できる\cite{vaswani2017attention,matsuo2006,ieice_shiori}．

%% =====================================================================
%% 第3章 提案手法
%% =====================================================================
\chapter{提案手法}
\label{chap:method}

\section{Xxxの手法}

% ---- 図と数式の例（不要になったら消す）----
提案手法の全体構成を図\ref{fig:overview}に示す．

\begin{figure}[tb]
  \centering
  \includegraphics[width=0.6\linewidth]{example-image-a} % ← 自分の図のファイル名にする
  \caption{提案手法の全体構成}
  \label{fig:overview}
\end{figure}

全結合層の出力 $y$ は式(\ref{eq:fc})で表される．
\begin{equation}
  y = f(Wx + b)
  \label{eq:fc}
\end{equation}
ここで，$W$は重み行列，$b$はバイアス，$f$は活性化関数である．

%% =====================================================================
%% 第4章 実験と結果
%% =====================================================================
\chapter{実験と結果}
\label{chap:experiment}

% ---- 表の例（不要になったら消す）----
実験に用いたFPGAボードの仕様を表\ref{tab:board}に示す．

\begin{table}[tb]
  \centering
  \caption{実験に用いたFPGAボードの仕様}
  \label{tab:board}
  \begin{tabular}{lr}
    \toprule
    項目 & 値 \\
    \midrule
    LUT & xxx,xxx \\
    DSP & x,xxx \\
    動作周波数 [MHz] & xxx \\
    \bottomrule
  \end{tabular}
\end{table}

%% =====================================================================
%% 第5章 考察と今後の課題
%% =====================================================================
\chapter{考察と今後の課題}
\label{chap:discussion}

%% =====================================================================
%% 第6章 結論
%% =====================================================================
\chapter{結論}
\label{chap:conclusion}

%% =====================================================================
%% 謝辞
%% =====================================================================
\chapter*{謝辞}
\addcontentsline{toc}{chapter}{謝辞}

ああああ

%% =====================================================================
%% 参考文献（references.bib から自動で作られる）
%% =====================================================================
\bibliographystyle{jIEEEtran}
\bibliography{references}

%% =====================================================================
%% 付録（無ければ \appendix から \end{document} の前までを消す）
%%   \appendix の後の \chapter は「付録A」「付録B」…になる．
%% =====================================================================
\appendix
\chapter{ソースコード}

付録として記載する事柄がある場合に限り使用する．

\begin{verbatim}
def infer(model, x):
    model.eval()
    return model(x)
\end{verbatim}

\end{document}
